% =========================================================
%  sections/01-identitas.tex
%  Bagian 1 — Identitas Mata Kuliah
% =========================================================

\section{Identitas Mata Kuliah}

\begin{table}[htbp]
\centering
\caption{Identitas mata kuliah \NamaMK}
\label{tab:identitas}
\begin{tabularx}{\textwidth}{@{}L{5.2cm}Y@{}}
\toprule
\headrow \textbf{Komponen} & \textbf{Keterangan} \\
\midrule
Nama Perguruan Tinggi   & \NamaPT \\
Nama Fakultas           & \NamaFakultas \\
Nama Program Studi      & \NamaProdi \\
Nama Mata Kuliah        & \NamaMK \\
Kode Mata Kuliah        & \KodeMK \\
Semester                & \SemesterMK \\
Bobot SKS               & 3 SKS (2 SKS teori, 1 SKS praktikum) \\
Sifat Mata Kuliah       & Wajib \\
Mata Kuliah Prasyarat   & Algoritma dan Pemrograman \\
Dosen Pengampu          & \NamaDosen \\
Bahasa Pengantar        & Bahasa Indonesia \\
Bentuk Pembelajaran     & Kuliah tatap muka, praktikum laboratorium komputer, dan tugas terstruktur \\
Media Pembelajaran      & LCD proyektor, komputer laboratorium, Python 3 + Pygame, Unity (versi LTS), repositori Git \\
Alokasi Waktu           & 16 minggu $\times$ 150 menit tatap muka (14 minggu materi, UTS, dan UAS) \\
Tahun Akademik          & \TahunAkademik \\
\bottomrule
\end{tabularx}
\end{table}

\subsection*{Deskripsi Singkat Mata Kuliah}

Mata kuliah \NamaMK\ membekali mahasiswa dengan konsep dan keterampilan dasar
pengembangan permainan digital. Perkuliahan disusun dalam dua tahap yang
berurutan.

\begin{enumerate}
  \item \textbf{Tahap fondasi algoritma (Python + Pygame).}
        Mahasiswa membangun pemahaman tentang \istilah{game} sebagai simulasi
        interaktif \istilah{real-time}, merancang \istilah{game loop} dan
        manajemen waktu, menggerakkan objek berbasis vektor dan
        \istilah{delta time}, menerapkan fisika sederhana, mendeteksi
        \istilah{collision}, serta mengimplementasikan kecerdasan buatan (AI)
        dasar untuk NPC berupa perilaku \istilah{chase}, \istilah{Finite State
        Machine} (FSM), dan \istilah{pathfinding} A*.
  \item \textbf{Tahap implementasi \istilah{game engine} (Unity).}
        Mahasiswa memindahkan solusi \istilah{computing} tersebut ke dalam
        Unity untuk membangun level, mengonfigurasi sistem fisika dan
        \istilah{collider}, serta mengintegrasikan aset visual dan audio ke
        dalam interaksi permainan.
\end{enumerate}

Pembelajaran berpusat pada mahasiswa melalui praktikum terbimbing, dua proyek
(mini proyek AI dan proyek akhir berupa prototipe \istilah{game}), serta
\istilah{playtesting} antar kelompok. Penilaian dilakukan berbasis kriteria
terhadap ketercapaian Sub-CPMK yang telah ditetapkan.

\subsection*{Bahan Kajian Utama}

\begin{itemize}
  \item Arsitektur \istilah{game loop}: siklus input--update--render dan manajemen waktu.
  \item Matematika vektor terapan, gerak berbasis \istilah{delta time}, dan integrasi numerik sederhana.
  \item \istilah{Collision detection} dan resolusi tabrakan pada lingkungan 2D.
  \item Kecerdasan buatan untuk NPC: perilaku reaktif, FSM, dan \istilah{pathfinding} A*.
  \item \istilah{Game engine} Unity: \istilah{scene}, \istilah{GameObject}, sistem fisika, dan \istilah{level design}.
  \item Integrasi aset visual, animasi, antarmuka pengguna, dan audio ke dalam \istilah{gameplay}.
  \item Evaluasi produk melalui \istilah{playtesting}, pelaporan \istilah{bug}, dan refleksi desain.
\end{itemize}
